
\documentclass[11pt,a4paper]{article}
\usepackage[a4paper,margin=2.05cm]{geometry}
\usepackage[spanish,es-noquoting]{babel}
\usepackage[T1]{fontenc}
\usepackage[utf8]{inputenc}
\usepackage{lmodern,microtype}
\usepackage{amsmath,amssymb,amsthm,mathtools}
\usepackage{booktabs,longtable,array,enumitem}
\usepackage{xcolor,hyperref}
\pagecolor{white}\color{black}
\setlength{\parindent}{0pt}
\setlength{\parskip}{0.65em}
\title{\bfseries N73 --- Teorema de puerta de orientación\\
\large \(\varphi\) fijo, \(\pi+e\) fijo, y supervivencia que selecciona la orientación}
\author{HMT / auditoría técnica}
\date{}
\begin{document}
\maketitle

\begin{abstract}
N73 formaliza la intuición de la novena puerta. La ambigüedad local de las fronteras profundas no es genérica: es una ambigüedad de orientación. La razón exacta es que \(q\) y \(a\) fijan completamente \(u^\varphi\) y \(u^\pi+u^e\). Por tanto, si varias ramas sobreviven localmente, sólo pueden diferir en cómo se reparte el par \(\pi/e\) alrededor de un eje fijo \(\varphi\). La auditoría \(t=5,\ldots,25\) verifica que todas las fronteras ambiguas tienen ese patrón, y que la supervivencia coinductiva resuelve la orientación en profundidad menor o igual que \(3\).
\end{abstract}

\section{Identidad exacta}

Para una frontera:
\[
B=
\begin{pmatrix}
u^\pi\\u^e\\u^\varphi
\end{pmatrix},
\]
las cantidades:
\[
q=u^\pi+u^e+u^\varphi,
\qquad
a=u^\pi+u^e-u^\varphi
\]
determinan:
\[
\boxed{
u^\varphi = 2(q-a).
}
\]
Y entonces:
\[
\boxed{
u^\pi+u^e = q-u^\varphi.
}
\]
Por tanto:
\[
\boxed{
\text{si }q,a\text{ están fijos, la única ambigüedad posible es el reparto de }\pi/e.
}
\]

\section{Auditoría de puertas ambiguas}

\[
\begin{array}{c|c|c|c|c}
t&\#S_t&u^\varphi&u^\pi+u^e&h_{\rm res}\\
\hline
5 & 2 & 221021 & 012100 & 2 \\
6 & 6 & 200221 & 101001 & 3 \\
7 & 11 & 110110 & 112000 & 3 \\
8 & 3 & 010122 & 120212 & 3 \\
9 & 3 & 112002 & 210110 & 2 \\
10 & 11 & 100021 & 102210 & 2 \\
11 & 3 & 000120 & 211201 & 2 \\
12 & 5 & 211122 & 220012 & 2 \\
13 & 3 & 202200 & 220021 & 2 \\
14 & 3 & 202110 & 222210 & 2 \\
15 & 5 & 221000 & 012212 & 2 \\
16 & 11 & 100102 & 012111 & 2 \\
17 & 6 & 111122 & 020122 & 2 \\
20 & 5 & 211021 & 211122 & 3 \\
21 & 2 & 101021 & 202011 & 3 \\
23 & 3 & 210002 & 222220 & 3 \\
24 & 3 & 012220 & 222121 & 2 \\
25 & 6 & 121011 & 112212 & 3 \\
\end{array}
\]

En todos los \(t\) ambiguos auditados:
\[
\boxed{
u^\varphi\text{ es constante a través de todas las ramas seleccionadas,}
}
\]
\[
\boxed{
u^\pi+u^e\text{ es constante a través de todas las ramas seleccionadas.}
}
\]
Así, las puertas ambiguas son puertas de orientación.

\section{Novena zona}

En la novena zona:
\[
t=8,\quad t=9,
\]
la pauta aparece con claridad:
\[
\varphi\text{ queda fijo,}
\qquad
\pi/e\text{ forman el espejo móvil,}
\]
y:
\[
\text{el futuro rompe el espejo.}
\]
Sin embargo, la auditoría muestra que esa pauta no es exclusiva de la novena puerta. La novena la hace visible; no la crea.

\section{Dictamen}

\[
\boxed{
\text{verde: la ambigüedad local profunda es ambigüedad de orientación.}
}
\]
\[
\boxed{
\text{verde: }\varphi\text{ actúa como eje fijo y }\pi/e\text{ como par móvil.}
}
\]
\[
\boxed{
\text{verde: la supervivencia coinductiva resuelve la orientación en profundidad }\le3.
}
\]
\[
\boxed{
\text{no se afirma aquí generación infinita completa.}
}
\]
\end{document}
